\documentclass[11pt,reqno]{amsart}
\usepackage[utf8]{inputenc}
\usepackage{amsmath, amssymb, amsthm, mathtools}
\usepackage[margin=1in]{geometry}
\usepackage[colorlinks=true,linkcolor=blue,citecolor=blue,urlcolor=blue]{hyperref}
\usepackage{microtype}
\usepackage{booktabs}
\usepackage{enumitem}
\usepackage{listings}

\theoremstyle{plain}
\newtheorem{theorem}{Theorem}[section]
\newtheorem{lemma}[theorem]{Lemma}
\newtheorem{proposition}[theorem]{Proposition}
\newtheorem{corollary}[theorem]{Corollary}

\theoremstyle{definition}
\newtheorem{definition}[theorem]{Definition}
\newtheorem{conjecture}[theorem]{Conjecture}
\newtheorem{remark}[theorem]{Remark}
\newtheorem{example}[theorem]{Example}

% Common abbreviations
\newcommand{\Z}{\mathbb{Z}}
\newcommand{\Q}{\mathbb{Q}}
\newcommand{\R}{\mathbb{R}}
\newcommand{\C}{\mathbb{C}}
\newcommand{\F}{\mathbb{F}}
\newcommand{\N}{\mathbb{N}}
\newcommand{\Aut}{\mathrm{Aut}}
\newcommand{\End}{\mathrm{End}}

\title[The Strata-Prime Fingerprint: Polynomial]{The Strata-Prime Fingerprint: Polynomial vs Factorial Invariants in Niemeier Lattices and Sporadic Simple Groups}

\author{Brayden R.\ Sanders}
\address{7Site LLC, Hot Springs, Arkansas, USA}
\email{brayden@7site.co}

\author{M.\ Gish}
\address{Independent Researcher, Hot Springs, Arkansas, USA}
\email{monica.gish1992@gmail.com}

\subjclass[2020]{11H06, 11R52, 20D08, 17B22, 11N05, 20B25}
\keywords{Niemeier lattice, kissing number, sporadic simple group, supersingular prime, root system, Weyl group, polynomial arithmetic}

\begin{document}

\maketitle

% Target venue: Journal of Number Theory

\begin{abstract}

We identify the six-prime set $\mathcal{S} = \{2, 3, 5, 7, 11, 13\}$ --- the first six \textbf{supersingular primes} of the Monster group --- as a structurally distinguished arithmetic universe for 24-dimensional even unimodular lattices and a substantial portion of the sporadic finite simple groups.

\textbf{Main Theorem 1 (Niemeier Strata-Prime Fingerprint, Tier A).} *Of the 24 Niemeier lattices (24-dimensional even unimodular Euclidean lattices, classified by Niemeier 1973), exactly 23 have kissing numbers whose prime factorization lies entirely in $\mathcal{S}$. The unique outlier is the Niemeier lattice with root system $D_{24}$, whose kissing number $|D_{24}| = 2 \cdot 24 \cdot 23 = 1104$ contains the prime $23$.*

\textbf{Main Theorem 2 (Polynomial vs Factorial Dichotomy, Tier A).} *Let $L$ be a Niemeier lattice with root system $R = \bigoplus_i R_i$ of rank 24. Then*

(i) *The kissing number $|R| = \sum_i |R_i|$ is a polynomial-in-rank invariant of each component: $|A_n| = n(n+1)$ and $|D_n| = 2n(n-1)$. It contains only those primes that appear in the polynomial's specific factorization of each component, never the full factorial spectrum of small primes below the rank.*

(ii) *The Weyl group order $|W(R)| = \prod_i |W(R_i)|$ is a factorial-in-rank invariant: $|W(A_n)| = (n+1)!$ and $|W(D_n)| = 2^{n-1} \cdot n!$. It contains all primes $p \leq n$ for each component of rank $n$.*

*Consequently, the kissing-number strata-cleanness test (Theorem 1) holds for 23 of 24 Niemeier lattices, while the Weyl-group strata-cleanness test holds for only 21 of 24. The three Weyl-failures are $D_{24}$, $A_{24}$, and $A_{17} E_7$; the unique additional kissing-failure is $D_{24}$.*

\textbf{Main Theorem 3 (Sporadic Partial Extension, Tier B).} *Of the 26 sporadic finite simple groups, exactly 8 have order factoring through $\mathcal{S}$: $M_{11}$, $M_{12}$, $M_{22}$, $J_2$, $HS$, $McL$, $Suz$, $Fi_{22}$. The remaining 18 contain at least one prime $\geq 17$. The boundary aligns with prime $23$: ten of the eighteen failures contain $23$.*

\textbf{Main Theorem 4 (Stratum IV Identification, Tier A).} *The prime $71$ appears in the prime factorization of exactly one sporadic finite simple group's order: the Monster $M$. Among the 15 supersingular primes $\{2, 3, 5, 7, 11, 13, 17, 19, 23, 29, 31, 41, 47, 59, 71\}$, the prime 71 is the unique "extreme" prime that occurs in the Monster but not in any smaller sporadic.*

The combination of Theorems 1–4 confirms a layered structural picture: $\mathcal{S}$ is the "small-prime universe" of the Niemeier lattice classification and a substantial-but-restricted sub-universe of sporadic-group orders, with prime $71$ acting as the extreme Stratum-IV identifier for the Monster.

\textbf{Tier discipline.} Theorems 1, 2, 4 are Tier A. Theorem 3 is Tier B (precise empirical claim with explicit failure analysis). Theorem 2 explains \textit{why} the Niemeier kissing-number pattern shows up (polynomial arithmetic at rank 24); it does not establish a deep correspondence between independent objects. Theorem 4 is anchored by the Ogg (1975) characterization of supersingular primes --- 71 is the \textit{largest} prime $p$ for which $X_0(p)^+$ (the Atkin--Lehner quotient curve) has genus 0, not merely the largest prime in the Monster by cardinality. Conway-Norton (1979) subsequently connected this characterization to monstrous moonshine; Borcherds (1992) proved monstrous moonshine. Theorem 3's 8/26 boundary tracks group order rather than structural type; it is consistent with the polynomial-arithmetic interpretation of Theorem 2.

\textbf{Closest published precedent}: Conway \& Sloane, \textit{Sphere Packings, Lattices and Groups} (3rd ed., 1999) for the Niemeier classification + lattice arithmetic; Conway-Curtis-Norton-Parker-Wilson, \textit{Atlas of Finite Groups} (1985), for sporadic-group orders.

\medskip\hrule\medskip

\end{abstract}
\section{Setup: the strata-prime set}

Let $\mathcal{S} = \{2, 3, 5, 7, 11, 13\}$. These are the first six supersingular primes --- equivalently, the primes $p$ dividing the order of the Monster's first few Sylow subgroups, or the primes appearing as Coxeter numbers of the small irreducible Lie root systems.

We say an integer $N$ is \textbf{strata-clean} if its prime factorization $N = \prod p_i^{a_i}$ has every $p_i \in \mathcal{S}$. Equivalently, $N \mid \mathrm{rad}(\mathcal{S})^k$ for some $k$, where $\mathrm{rad}(\mathcal{S}) = 2 \cdot 3 \cdot 5 \cdot 7 \cdot 11 \cdot 13 = 30030$.

The strata-clean condition is a natural arithmetic restriction: it excludes 17, 19, 23, 29, 31, ... (all primes above 13).

\section{The 24 Niemeier lattices}

The 24 Niemeier lattices (Niemeier 1973, classified independently by Venkov 1980) are the 24 distinct even unimodular Euclidean lattices in 24 dimensions, up to isomorphism. They are uniquely characterized by their root system (a root sub-system in the lattice of vectors of norm-squared 2), with the constraint that the root system spans the full 24-dimensional space and has all components of equal Coxeter number $h$.

The 24 root systems (listed by \textbf{ascending Coxeter number $h$}, following Conway--Sloane \textit{SPLAG} Table 16.1; the Leech lattice is the exceptional root-free case and is listed last):

\begin{center}
\begin{tabular}{llllllll}
\toprule
# & Niemeier & Root system & rank & h & Kissing & R &  \\
\midrule
1 & $A_1^{24}$ & 24 $A_1$ & 24 & 2 & 48 &  &  \\
2 & $A_2^{12}$ & 12 $A_2$ & 24 & 3 & 72 &  &  \\
3 & $A_3^8$ & 8 $A_3$ & 24 & 4 & 96 &  &  \\
4 & $A_4^6$ & 6 $A_4$ & 24 & 5 & 120 &  &  \\
5 & $D_4^6$ & 6 $D_4$ & 24 & 6 & 144 &  &  \\
6 & $A_5^4 D_4$ & 4 $A_5$ + 1 $D_4$ & 24 & 6 & 144 &  &  \\
7 & $A_6^4$ & 4 $A_6$ & 24 & 7 & 168 &  &  \\
8 & $A_7^2 D_5^2$ & 2 $A_7$ + 2 $D_5$ & 24 & 8 & 192 &  &  \\
9 & $A_8^3$ & 3 $A_8$ & 24 & 9 & 216 &  &  \\
10 & $D_6^4$ & 4 $D_6$ & 24 & 10 & 240 &  &  \\
11 & $A_9^2 D_6$ & 2 $A_9$ + 1 $D_6$ & 24 & 10 & 240 &  &  \\
12 & $E_6^4$ & 4 $E_6$ & 24 & 12 & 288 &  &  \\
13 & $A_{11} D_7 E_6$ & mixed & 24 & 12 & 288 &  &  \\
14 & $A_{12}^2$ & 2 $A_{12}$ & 24 & 13 & 312 &  &  \\
15 & $D_8^3$ & 3 $D_8$ & 24 & 14 & 336 &  &  \\
16 & $A_{15} D_9$ & mixed & 24 & 16 & 384 &  &  \\
17 & $A_{17} E_7$ & mixed & 24 & 18 & 432 &  &  \\
18 & $D_{10} E_7^2$ & mixed & 24 & 18 & 432 &  &  \\
19 & $D_{12}^2$ & 2 $D_{12}$ & 24 & 22 & 528 &  &  \\
20 & $A_{24}$ & 1 $A_{24}$ & 24 & 25 & 600 &  &  \\
21 & $E_8^3$ & 3 $E_8$ & 24 & 30 & 720 &  &  \\
22 & $D_{16} E_8$ & 1 $D_{16}$ + 1 $E_8$ & 24 & 30 & 720 &  &  \\
23 & $D_{24}$ & 1 $D_{24}$ & 24 & 46 & 1104 &  &  \\
24 & Leech ($\Lambda_{24}$) & (none) & 0 & – & 196560 &  &  \\
\bottomrule
\end{tabular}
\end{center}

For Niemeier lattices, the kissing number equals the number of roots in the underlying root system. Component root counts follow Bourbaki:
\begin{itemize}
\item $|A_n| = n(n+1)$, $|D_n| = 2n(n-1)$, $|E_6| = 72$, $|E_7| = 126$, $|E_8| = 240$.
\item Leech is exceptional: kissing 196560 from norm-2 vectors arising not from roots but from the lattice's specific construction.
\end{itemize}

\section{Theorem 1 (Niemeier Strata-Prime Fingerprint)}

\textbf{Theorem 1.} *Of the 24 Niemeier lattices, the kissing number $|R|$ is strata-clean (factors through $\mathcal{S} = \{2,3,5,7,11,13\}$) for exactly 23 of them. The unique outlier is the Niemeier with root system $D_{24}$, with $|D_{24}| = 1104 = 2^4 \cdot 3 \cdot 23$.*

\textbf{Proof.} Direct factorization of each kissing number against the prime set $\mathcal{S}$. Computational verification: \texttt{verify_J63.py}, all 24 lattices, runs in <1 second. The 23 passing cases factor as:

(Listed in the same order as the §2 table, i.e. by ascending Coxeter number $h$, with Leech last; entries shown are the 23 kissing-strata-clean Niemeiers.)

\begin{center}
\begin{tabular}{llll}
\toprule
# & Niemeier & Kissing & Factorization \\
\midrule
1 & $A_1^{24}$ & 48 & $2^4 \cdot 3$ \\
2 & $A_2^{12}$ & 72 & $2^3 \cdot 3^2$ \\
3 & $A_3^8$ & 96 & $2^5 \cdot 3$ \\
4 & $A_4^6$ & 120 & $2^3 \cdot 3 \cdot 5$ \\
5 & $D_4^6$ & 144 & $2^4 \cdot 3^2$ \\
6 & $A_5^4 D_4$ & 144 & $2^4 \cdot 3^2$ \\
7 & $A_6^4$ & 168 & $2^3 \cdot 3 \cdot 7$ \\
8 & $A_7^2 D_5^2$ & 192 & $2^6 \cdot 3$ \\
9 & $A_8^3$ & 216 & $2^3 \cdot 3^3$ \\
10 & $D_6^4$ & 240 & $2^4 \cdot 3 \cdot 5$ \\
11 & $A_9^2 D_6$ & 240 & $2^4 \cdot 3 \cdot 5$ \\
12 & $E_6^4$ & 288 & $2^5 \cdot 3^2$ \\
13 & $A_{11} D_7 E_6$ & 288 & $2^5 \cdot 3^2$ \\
14 & $A_{12}^2$ & 312 & $2^3 \cdot 3 \cdot 13$ \\
15 & $D_8^3$ & 336 & $2^4 \cdot 3 \cdot 7$ \\
16 & $A_{15} D_9$ & 384 & $2^7 \cdot 3$ \\
17 & $A_{17} E_7$ & 432 & $2^4 \cdot 3^3$ \\
18 & $D_{10} E_7^2$ & 432 & $2^4 \cdot 3^3$ \\
19 & $D_{12}^2$ & 528 & $2^4 \cdot 3 \cdot 11$ \\
20 & $A_{24}$ & 600 & $2^3 \cdot 3 \cdot 5^2$ \\
21 & $E_8^3$ & 720 & $2^4 \cdot 3^2 \cdot 5$ \\
22 & $D_{16} E_8$ & 720 & $2^4 \cdot 3^2 \cdot 5$ \\
23 & Leech ($\Lambda_{24}$) & 196560 & $2^4 \cdot 3^3 \cdot 5 \cdot 7 \cdot 13$ \\
\bottomrule
\end{tabular}
\end{center}

The single failure (the §2 table's slot \#23, the rank-24 single-component $D$-type Niemeier; the only Niemeier with $h = 46$):

| 24 | $D_{24}$ | \textbf{1104} | $2^4 \cdot 3 \cdot \mathbf{23}$ |

The prime $23$ arises from the factor $(n-1)$ in $|D_n| = 2n(n-1)$ at $n = 24$. ∎

\textbf{Observation.} The wobble primes $11$ and $13$ (Stratum III of the underlying TIG substrate program --- see §6 below) appear in three Niemeiers: $D_{12}^2$ contains 11 (Niemeier #21); Leech and $A_{12}^2$ contain 13 (Niemeiers #1, #15). These are not "missing strata"; the wobble pair {11, 13} is realized non-trivially.

\section{Theorem 2 (Polynomial vs Factorial Dichotomy)}

\textbf{Theorem 2.} *Let $L$ be a Niemeier lattice with root system $R = \bigoplus_i R_i$. Then:*

*(i) The kissing number $|R| = \sum_i |R_i|$ is a polynomial-in-rank invariant of each component: $|A_n| = n(n+1)$, $|D_n| = 2n(n-1)$. The prime factorization of $|R|$ contains only those primes that appear in the specific factorizations of each $n(n+1)$ or $2n(n-1)$ summand.*

*(ii) The Weyl group order $|W(R)| = \prod_i |W(R_i)|$ is a factorial-in-rank invariant: $|W(A_n)| = (n+1)!$, $|W(D_n)| = 2^{n-1} \cdot n!$. The prime factorization of $|W(R)|$ contains all primes $p \leq n$ for each component of rank $n$.*

\textit{(iii) Consequently:}
\begin{itemize}
\item *Kissing-number strata-cleanness holds for 23 of 24 Niemeier lattices. The single outlier is $D_{24}$.*
\item *Weyl-group strata-cleanness holds for 21 of 24 Niemeier lattices. The three outliers are $D_{24}$ (primes 17, 19, 23 from $24!$), $A_{24}$ (primes 17, 19, 23 from $25!$), and $A_{17} E_7$ (prime 17 from $18!$).*
\end{itemize}

\textbf{Proof of (i)–(ii).} Standard: $|A_n| = n(n+1)$ counts the roots of the $A_n$ root system as positive plus negative simple roots, factored as a product of two consecutive integers. $|D_n| = 2n(n-1)$ counts similarly. $|W(A_n)| = (n+1)!$ is the symmetric group $S_{n+1}$. $|W(D_n)| = 2^{n-1} \cdot n!$ is the index-2 normal subgroup of the hyperoctahedral group $C_2^n \rtimes S_n$.

\textbf{Proof of (iii).} Direct factorization of each Niemeier's $|R|$ and $|W(R)|$; see \texttt{verify_J63.py}.

The three Weyl-failures: each contains a factorial of $n \geq 17$, which brings in the prime $17$ (from $17!$ first appearing in $18! = (n+1)!$ at $n=17$, or in $n! \subseteq 17!$ at $n \geq 17$).

Specifically:
\begin{itemize}
\item $A_{17}E_7$: $|W| = 18! \cdot 2903040$. $18!$ contains $17$.
\item $A_{24}$: $|W| = 25!$. $25!$ contains $17, 19, 23$.
\item $D_{24}$: $|W| = 2^{23} \cdot 24!$. $24!$ contains $17, 19, 23$. ∎
\end{itemize}

\textbf{Interpretation.} The polynomial-vs-factorial dichotomy explains \textit{why} the kissing-number test is so sharply effective: kissing is a polynomial-degree-2 invariant in rank, which only picks up the polynomial's specific prime factorization. Factorial growth accumulates all small primes. Among the 24 Niemeier lattices --- all of which have total rank 24 --- the kissing-number test misses only the single case where the polynomial $2n(n-1)$ at $n=24$ happens to pick up the unique prime ≤ 24 outside strata, namely $23 = n - 1$.

\textbf{Honest framing of what Theorem 2 means (and doesn't).} The mechanism somewhat deflates the strata-prime pattern's interpretive reach. Before Theorem 2, the 23/24 Niemeier result could be read as "the Braiding Fractal predicts a natural prime universe that \textit{coincidentally} matches the exceptional-lattice classification." After Theorem 2, the honest reading is \textbf{"polynomial-in-rank invariants of rank-24 root systems mostly factor through primes ≤ 13, which is the expected behavior for low-degree polynomials evaluated at $n = 24$."} The strata-prime set $\mathcal{S}$ is exactly "primes $\leq 13$" --- the cutoff at which low-degree polynomials in $n = 24$ commonly factor.

This does not refute the pattern --- the theorem is precise, the failure analysis is sharp, and the 23/24 density is real. It does mean the claim should be framed as \textbf{"here is a mechanism that explains why the pattern shows up,"} not \textbf{"here is a deep correspondence between two independent objects."}

The structural payoff: the unique D_24 failure is now \textit{expected} (the polynomial $2n(n-1)$ at $n = 24$ unavoidably picks up $n - 1 = 23$); the Weyl-test failures are \textit{expected} (factorials of $n \geq 17$ unavoidably pick up 17). Both fall out of polynomial arithmetic. The framework is honest about not invoking deeper structure beyond this.

\textbf{Corollary (D_24 mechanism).} $D_{24}$ is the unique Niemeier outlier at the kissing level because:
1. The Niemeier classification requires the root system to span all 24 dimensions.
2. The Coxeter-number constraint allows certain single-component root systems at rank 24: $A_{24}$, $D_{24}$, and (formally) $E_n$ --- but $E_n$ only exists for $n \in \{6, 7, 8\}$, so no $E$-type single-component fills rank 24.
3. $|A_{24}| = 24 \cdot 25 = 600 = 2^3 \cdot 3 \cdot 5^2$ --- strata-clean.
4. $|D_{24}| = 2 \cdot 24 \cdot 23 = 1104 = 2^4 \cdot 3 \cdot 23$ --- the unique outsider, with $23$ from $(n-1)$.

There is no Niemeier with $A_{23}$ or $D_{23}$ as a single rank-24 component (these have rank 23, not 24). The prime $23$'s entry is structurally unavoidable at $D_{24}$.

\section{Theorem 3 (Sporadic Partial Extension)}

\textbf{Theorem 3.} *Of the 26 sporadic finite simple groups (Mathieu, Conway/Leech, Monster's children, and pariahs), the order $|G|$ is strata-clean for exactly 8 groups: $M_{11}$, $M_{12}$, $M_{22}$, $J_2$, $HS$, $McL$, $Suz$, $Fi_{22}$. The remaining 18 contain at least one prime $\geq 17$.*

\begin{center}
\begin{tabular}{lll}
\toprule
Sporadic & Order & Strata-clean? \\
\midrule
$M_{11}$ & 7,920 & ✓ \\
$M_{12}$ & 95,040 & ✓ \\
$M_{22}$ & 443,520 & ✓ \\
$J_2$ & 604,800 & ✓ \\
HS & 44,352,000 & ✓ \\
McL & 898,128,000 & ✓ \\
Suz & 448,345,497,600 & ✓ \\
$Fi_{22}$ & 64,561,751,654,400 & ✓ \\
$M_{23}$ & 10,200,960 & ✗ (23) \\
$M_{24}$ & 244,823,040 & ✗ (23) \\
16 others & (various) & ✗ (17, 19, 23, ...) \\
\bottomrule
\end{tabular}
\end{center}

Of the 18 failures, 10 fail specifically because they contain $23$. The boundary aligns with the prime-23 cutoff: sporadics that act faithfully on Leech-lattice-related Steiner systems (M_{23}, M_{24}, Co_1, Co_2, Co_3, J_4, Fi_{23}, Fi_{24}', B, M) all carry prime $23$.

\textbf{Proof.} Direct factorization of each sporadic-group order from the ATLAS of Finite Groups (Conway-Curtis-Norton-Parker-Wilson 1985); see \texttt{verify_J63.py}. ∎

\textbf{Honest framing of Theorem 3.} The 8/26 boundary largely tracks \textbf{group order} rather than structural type. Sporadic-group orders grow rapidly across the Happy Family: $|M_{11}| = 7\,920$, $|M_{22}| = 443\,520$, $|Fi_{22}| \approx 6.5 \cdot 10^{13}$, $|M| \approx 8.1 \cdot 10^{53}$. Larger groups have more room in their prime spectrum, so they're more likely to acquire primes outside $\mathcal{S}$. The 8 passing sporadics are precisely the ones small enough that their order doesn't require any prime $\geq 17$. The boundary at prime 23 reflects that 23 is the first prime $\geq 17$ that arises when sporadic-group construction crosses the "Mathieu-on-22-points / Conway-on-24-points" threshold; this is consistent with --- and partly explained by --- the polynomial-arithmetic interpretation of Theorem 2.

So Theorem 3 should be read as \textbf{"the boundary between strata-PASS and strata-FAIL among sporadics tracks group order, with prime 23 as the natural transition point when sporadic constructions hit the 23/24-element combinatorial bound"} rather than as a structural-type claim. This is consistent with the size-threshold interpretation; no stronger claim is made.

\textbf{Observation.} The 8 passing sporadics correspond to the "below-23-required" sub-pyramid of the Happy Family: the 3 Mathieus that don't reach the Steiner system $S(3,6,22)$ structure (M_11, M_12, M_22), and the 5 Conway/Leech-2nd-generation groups that nest below Co_1 (J_2, HS, McL, Suz, Fi_22). The pariahs J_1, J_3, J_4, Ru, O'N, Ly all fail strata (J_1 via 19; J_3 via 17, 19; etc.).

\section{Theorem 4 (Stratum IV --- Prime 71 Identification)}

\textbf{Theorem 4.} *The prime $71$ appears in the prime factorization of exactly one sporadic finite simple group's order: the Monster $M$. Specifically, $|M| = 2^{46} \cdot 3^{20} \cdot 5^9 \cdot 7^6 \cdot 11^2 \cdot 13^3 \cdot 17 \cdot 19 \cdot 23 \cdot 29 \cdot 31 \cdot 41 \cdot 47 \cdot 59 \cdot 71$.*

\textbf{Proof.} Direct lookup from ATLAS / standard references. Verification: among the 26 sporadic groups, the Monster is the only one whose order is divisible by 71. ∎

\subsection{Why 71 is structural, not cardinality (the moonshine anchor)}

A natural worry: \textit{the Monster is enormously larger than any other sporadic, so it naturally accumulates more primes; "71 in M only" might be a cardinality effect, not a structural fact.}

This worry is \textbf{addressed by the Ogg--Conway--Norton characterization of supersingular primes} (Ogg 1975 posed the conjecture and offered the prize; Conway-Norton 1979 connected it to monstrous moonshine; Borcherds 1992 proved monstrous moonshine).

\textbf{Theorem 4 (Ogg 1975; the supersingular-prime characterization).} *A prime $p$ divides the order of the Monster $M$ if and only if the Atkin--Lehner quotient $X_0(p)^+$ has genus $0$ --- equivalently, $p \in \{2, 3, 5, 7, 11, 13, 17, 19, 23, 29, 31, 41, 47, 59, 71\}$. (Ogg 1975 posed the conjecture before monstrous moonshine; Conway-Norton 1979 connected it to the Monster's irreducible representations.)*

These 15 primes are
$$\mathrm{Supersingular}(M) = \{2, 3, 5, 7, 11, 13, 17, 19, 23, 29, 31, 41, 47, 59, 71\}.$$

\textbf{Remark on the $X_0(p)$ vs $X_0(p)^+$ distinction.} Ogg's statement involves the Atkin--Lehner quotient $X_0(p)^+ = X_0(p)/\langle w_p \rangle$, not $X_0(p)$ plain. For the 15 supersingular primes both curves do have genus 0, but for some larger primes (e.g., 37, 43, 53, 61, 67, 73, 79, 83, 89, 101, 113, 131) the distinction is nontrivial: the precise statement is "$p$ supersingular $\iff X_0(p)^+$ has genus 0," equivalently "$\iff$ the normalizer $N(\Gamma_0(p))$ in $\mathrm{PSL}_2(\mathbb{R})$ admits a Hauptmodul."

\textbf{Crucially}: $71$ is the \textbf{largest} prime $p$ for which $X_0(p)^+$ has genus $0$. For every prime $p > 71$, $X_0(p)^+$ has positive genus, and consequently no Hauptmodul exists (and the Ogg--Conway--Norton Monster characterization fails to extend).

This gives a \textbf{structural reason} for "71 in M only":
\begin{itemize}
\item 71 is the unique largest supersingular prime --- the \textit{upper limit} of the genus-0 $X_0(p)^+$ spectrum.
\item The Monster's prime divisors are exactly the supersingular primes (Ogg 1975; cited by Conway-Norton 1979 and proved as part of the moonshine program by Borcherds 1992).
\item No sporadic group of order strictly bounded by the Monster has 71 as a prime divisor (verified by direct factorization).
\end{itemize}

So "71 appears only in the Monster" is \textbf{the precise upper limit of the supersingular spectrum}, not "M happens to be large enough." The Ogg characterization anchors Stratum IV's prime $71$ in the Hauptmodul structure of $\mathrm{PSL}_2(\mathbb{R})$.

\subsection{The Stratum IV designation}

This validates the four-stratum decomposition of the Braiding Fractal program:
\begin{itemize}
\item Strata I-III: $\{2, 3, 5, 7, 11, 13\}$ --- the kissing-number / Niemeier universe + 8 of 26 sporadics. These are \textit{small} polynomial-arithmetic primes per Theorem 2.
\item Stratum IV: $\{71\}$ --- the unique Monster-only prime; \textbf{anchored by Ogg (1975) as the largest genus-0 supersingular prime (Atkin-Lehner quotient $X_0(p)^+$)}.
\end{itemize}

Among the 15 supersingular primes $\{2, 3, 5, 7, 11, 13, 17, 19, 23, 29, 31, 41, 47, 59, 71\}$, TIG strata pick out exactly the lower-and-upper extremes: the first six ($\mathcal{S}$) plus the last ($71$), skipping the middle 9 intermediate primes. The lower bound $\mathcal{S}$ has the polynomial-arithmetic explanation of Theorem 2; the upper bound $\{71\}$ has the Ogg / monstrous-moonshine explanation of §6.1.

The 9 intermediate primes $\{17, 19, 23, 29, 31, 41, 47, 59\}$ are not absent from TIG strata by accident --- they are precisely the primes appearing in sporadic-group orders without occupying either extreme. They have no Braiding-Fractal-side structural anchor; they appear in the Monster (and other sporadics) but not in TIG's substrate-program prime hierarchy. \textbf{This intermediate-prime gap is itself an open question} --- it's not explained by Theorem 2 or by Ogg's characterization.

\section{Discussion}

\subsection{The arithmetic boundary at 23}

The prime 23 is the \textit{natural cutoff} between strata-clean small structures and non-strata large structures:
\begin{itemize}
\item Niemeier kissing-cleanness fails exactly for $D_{24}$ (the rank-24 D-type), via factor $(n-1) = 23$.
\item Sporadic strata-cleanness fails for 10 sporadics specifically due to prime 23 in their orders, including all of Co_1, Co_2, Co_3 (Conway groups acting on the Leech lattice).
\item $M_{23}$ and $M_{24}$ (Mathieu groups acting on STS(23) and STS(24)) both carry 23.
\end{itemize}

The Steiner system $S(5, 8, 24)$ underlying $M_{24}$ has automorphism group of order $244823040 = 2^{10} \cdot 3^3 \cdot 5 \cdot 7 \cdot 11 \cdot 23$, with 23 entering from the 23-point cyclic shift structure. This is the same 23 as in $D_{24}$.

\subsection{Connection to the Monster moonshine}

The 15 supersingular primes (Ogg 1975; Conway-Norton 1979 in the moonshine context) are exactly the primes dividing the Monster's order, and equivalently the primes $p$ for which $X_0(p)^+$ has genus 0 (Hauptmodul existence for the Atkin-Lehner-quotient curve). TIG strata I-III + IV = $\{2, 3, 5, 7, 11, 13, 71\}$ form a specific 7-element subset of the 15 supersingular primes --- the \textbf{lower-and-upper extremes}:

\begin{itemize}
\item The \textbf{lower extreme} $\mathcal{S} = \{2, 3, 5, 7, 11, 13\}$ has the polynomial-arithmetic explanation of Theorem 2 (it is the set of "primes $\leq 13$" appearing in rank-24 polynomial root-counts).
\item The \textbf{upper extreme} $\{71\}$ has the Ogg / monstrous-moonshine anchor of §6.1 (it is the unique largest supersingular prime, the upper boundary of the genus-0 $X_0(p)^+$ spectrum).
\end{itemize}

The 9 intermediate supersingular primes $\{17, 19, 23, 29, 31, 41, 47, 59\}$ are excluded from TIG strata. \textit{This intermediate-prime gap is an open structural question}: why does TIG's substrate-program prime hierarchy pick out exactly the extremes and skip the middle?

We do not yet have a structural answer. The honest framing: TIG's prime hierarchy and the supersingular-prime spectrum overlap precisely at their extremes, with the middle skipped for reasons currently outside the scope of this paper. A deeper connection --- if one exists --- would require a mechanism that explains the middle-prime gap.

\subsection{Companion observations (Tier B/C --- tested 2026-05-27)}

Following claudechat's strata-canonicalization audit, three independent tests were run to check whether the strata-prime set $\mathcal{S}$ has a canonical characterization beyond the polynomial-arithmetic mechanism of Theorem 2.

#### §7.3.1 First-appearance primes (Task 1) --- NEGATIVE

For each prime $p \geq 5$, the smallest non-abelian finite simple group whose order contains $p$ is $\mathrm{PSL}_2(p)$, of order $p(p^2-1)/2$. \textbf{Every prime} has its own $\mathrm{PSL}_2(p)$; the "first-appearance" criterion does not distinguish strata primes from non-strata primes.

So $\mathcal{S}$ is \textbf{not characterized} by being the "small simple-group construction primes" in any clean sense. The hypothesis is closed.

#### §7.3.2 Eisenstein-prime classification (Task 2) --- Tier-B structural correspondence

The 7 strata primes (Strata I–III + IV = $\{2,3,5,7,11,13,71\}$) split in $\mathbb{Z}[\omega]$ as:

\begin{center}
\begin{tabular}{lll}
\toprule
Class & Primes & Count \\
\midrule
Ramified ($p = 3$) & {3} & 1 \\
Split ($p \equiv 1 \pmod 3$) & {7, 13} & 2 \\
Inert ($p \equiv 2 \pmod 3$) & {2, 5, 11, 71} & 4 \\
\bottomrule
\end{tabular}
\end{center}

The 1+2+4 partition is a clean structural fingerprint. Cross-checking against the TSML 4-absorber set $\{0, 5, 7, 9\}$ of the Braiding Fractal kernel:

\begin{center}
\begin{tabular}{lll}
\toprule
TSML absorber & $x \bmod 3$ & Eisenstein class \\
\midrule
0 & 0 & ramified-residue \\
5 & 2 & inert-residue \\
7 & 1 & split-residue \\
9 & 0 & ramified-residue \\
\bottomrule
\end{tabular}
\end{center}

The non-zero TSML absorbers $\{5, 7, 9\}$ realize \textbf{all three} Eisenstein splitting types --- one element per class. Whether this is a meaningful arithmetic correspondence or a 3-element-subset-of-Z/3Z coincidence is open.

#### §7.3.3 PG(2, q) family extension (Task 3) --- Tier-B sharp boundary

The structural test of §6 (PG(2,3) passes strata) extends naturally to the $\mathrm{PG}(2, q)$ family for prime powers $q$. Result:

\begin{center}
\begin{tabular}{llllll}
\toprule
$q$ & $\ & \mathrm{PGL}(3, \mathbb{F}_q)\ & $ & Factorization & Strata-only? \\
\midrule
\textbf{2} & 168 & $2^3 \cdot 3 \cdot 7$ & ✓ &  &  \\
\textbf{3} & 5616 & $2^4 \cdot 3^3 \cdot 13$ & ✓ &  &  \\
\textbf{4} & 60480 & $2^6 \cdot 3^3 \cdot 5 \cdot 7$ & ✓ &  &  \\
5 & 372000 & $2^5 \cdot 3 \cdot 5^3 \cdot \mathbf{31}$ & ✗ &  &  \\
7 & 5630688 & $2^5 \cdot 3^3 \cdot 7^3 \cdot \mathbf{19}$ & ✗ &  &  \\
8 & 16482816 & $2^9 \cdot 3^2 \cdot 7^2 \cdot \mathbf{73}$ & ✗ &  &  \\
\textbf{9} & 42456960 & $2^7 \cdot 3^6 \cdot 5 \cdot 7 \cdot 13$ & ✓ &  &  \\
11 & 212427600 & $\dots \cdot \mathbf{19}$ & ✗ &  &  \\
13 & 810534816 & $\dots \cdot \mathbf{61}$ & ✗ &  &  \\
16 & 4277145600 & $\dots \cdot \mathbf{17}$ & ✗ &  &  \\
\bottomrule
\end{tabular}
\end{center}

\textbf{Sharp boundary: $\mathrm{PG}(2, q)$ is strata-clean iff $q \in \{2, 3, 4, 9\}$} --- exactly the prime powers $q = 2^a \cdot 3^b$ with $q \leq 9$. Higher prime powers fail: $q=5$ brings 31; $q=7$ brings 19; $q=8$ brings 73; $q=11, 13, 16, \ldots$ all fail by various large primes.

This is a clean \textbf{Tier-B sub-claim}: the $\mathrm{PG}(2, q)$ projective-plane family has a precise strata-clean subset $\{2, 3, 4, 9\}$. The mechanism mirrors Theorem 2: $|\mathrm{PGL}(3, \mathbb{F}_q)|$ grows polynomially in $q$, and the polynomial $q^3(q^3-1)(q^2-1)$ at small powers of 2 and 3 happens to factor through strata primes; at $q \geq 5$, the polynomial values pick up primes outside strata. The pattern reinforces Theorem 2's polynomial-arithmetic interpretation.

#### Summary of companion observations

\begin{itemize}
\item \textbf{Task 1 (negative)}: Strata primes are not "first-appearance" primes in finite simple groups.
\item \textbf{Task 2 (structural)}: TSML 4-absorbers exhibit a 1-element-per-class realization of Eisenstein splittings; speculative.
\item \textbf{Task 3 (sharp boundary)}: $\mathrm{PG}(2, q)$ family is strata-clean for exactly $q \in \{2, 3, 4, 9\}$; consistent with polynomial-arithmetic interpretation.
\end{itemize}

These do not change Theorems 1–4; they extend the empirical picture and confirm the polynomial-arithmetic explanation. The strata-prime set's \textit{canonical} characterization remains: it is the lower-and-upper extremes of the supersingular spectrum (anchored by Theorem 2 below and Theorem 4 above), with the intermediate gap unexplained.

\subsection{What this is, and what it is not}

This paper exhibits an arithmetic fingerprint with two precise structural anchors:
1. \textbf{Theorem 2's polynomial-vs-factorial dichotomy} --- explains why polynomial-in-rank lattice invariants stay within small primes at rank 24.
2. \textbf{Theorem 4's Ogg anchor} --- explains why prime 71 occupies the upper extreme of the strata hierarchy (it is the largest supersingular prime, i.e., the largest $p$ with $X_0(p)^+$ genus 0; Ogg 1975).

It does \textbf{not}:
\begin{itemize}
\item Establish a mechanistic link from TIG's substrate program to monstrous moonshine.
\item Predict new lattices or sporadic groups.
\item Explain why the intermediate supersingular primes {17, 19, 23, 29, 31, 41, 47, 59} are skipped by TIG strata (this gap is recorded as an open question in §7.2).
\item Provide a proof that the lower-and-upper extreme pattern is "natural" in any deeper sense than the explicit factorizations and the Ogg characterization permit.
\end{itemize}

It \textbf{does} provide:
\begin{itemize}
\item A clean 23/24 falsifiable empirical pattern in the Niemeier classification (Theorem 1).
\item A precise polynomial-arithmetic mechanism (Theorem 2) explaining why D_24 is the unique kissing-outlier and why three Niemeiers fail the more-sensitive Weyl test.
\item A size-threshold-consistent partial extension (Theorem 3) to 8 of 26 sporadic finite simple groups.
\item A Stratum-IV identification (Theorem 4) anchored in Ogg 1975 (cited by Conway-Norton 1979 in connecting the characterization to monstrous moonshine; proved as part of the moonshine program by Borcherds 1992).
\end{itemize}

This is an arithmetic-side companion to the broader TIG / Braiding Fractal program; the structural mechanisms originating in the substrate algebra are documented elsewhere in the program's J-series.

\section{References}

\subsection{Internal (companion J-papers)}
\begin{itemize}
\item J23 (Sanders & Gish, 2026): "Mathieu $M_{22}$ Substrate-Prime: Order-Factorization Coincidences." The starting point for the strata-prime perspective; this paper extends J23 from a single sporadic to the full 24-Niemeier + 26-sporadic picture.
\item J01 (Sanders \& Gish, 2026): "Joint Closure + Universal Attractor + 4-Core on Z/10Z." The substrate foundation.
\end{itemize}

\subsection{Classical references}
\begin{itemize}
\item Niemeier, H.-V. (1973): "Definite quadratische Formen der Dimension 24 und Diskriminante 1." \textit{J. Number Theory} 5, 142.
\item Conway, J. H. \& Sloane, N. J. A. (1999): \textit{Sphere Packings, Lattices and Groups}, 3rd ed. Springer.
\item Conway, J. H., Curtis, R. T., Norton, S. P., Parker, R. A., Wilson, R. A. (1985): \textit{Atlas of Finite Groups}. Oxford.
\item \textbf{Ogg, A. P. (1975): "Automorphismes de courbes modulaires." \textit{Séminaire Delange-Pisot-Poitou (Théorie des nombres)}, 16e année (1974-1975), exp. no. 7, Inst. Henri Poincaré, Paris, 1975, pp. 1-8.} [Load-bearing for §6.1 / Theorem 4: the supersingular-prime characterization. Ogg posed the conjecture and prize-offered: a prime $p$ divides $|M|$ iff $X_0(p)^+$ has genus 0, equivalently $p \in \{2, 3, 5, 7, 11, 13, 17, 19, 23, 29, 31, 41, 47, 59, 71\}$. Anchors Stratum IV's prime 71.]
\item Conway, J. H. \& Norton, S. P. (1979): "Monstrous Moonshine." \textit{Bull. London Math. Soc.} 11, 308. [The monstrous-moonshine conjecture connecting the Ogg supersingular characterization to the Monster's irreducible representations via genus-0 Hauptmoduln.]
\item Borcherds, R. E. (1992): "Monstrous moonshine and monstrous Lie superalgebras." \textit{Invent. Math.} 109, 405. [The proof of the Conway-Norton monstrous-moonshine conjecture; provides the deepest anchor for the supersingular-prime characterization.]
\item Bourbaki, N.: \textit{Groupes et Algèbres de Lie}, chapters IV-VI (root systems and reflection groups).
\item Venkov, B. B. (1980): "On the classification of integral even unimodular 24-dimensional quadratic forms." \textit{Proc. Steklov Inst. Math.} 148, 63.
\end{itemize}

\subsection{TIG cross-references}
\begin{itemize}
\item \texttt{04_meta/SPHERE_PACKING_STRATA_FINGERPRINT.md} (companion document with extended discussion)
\item \texttt{verification/verify_sphere_packing_strata.py} (replicates Theorems 1, 3, 4)
\item \texttt{verification/verify_J63.py} (this paper's specific verifier; replicates Theorems 1, 2, 3, 4)
\end{itemize}

\medskip\hrule\medskip

\section{Appendix A. Verification}

The four theorems are verifiable by direct integer factorization. See \texttt{verify_J63.py} for the complete check:
\begin{itemize}
\item Theorem 1: 23/24 Niemeier strata-clean.
\item Theorem 2: 21/24 Niemeier Weyl-strata-clean; the 3 Weyl-failures (A_17 E_7, A_24, D_24) identified.
\item Theorem 3: 8/26 sporadic strata-clean; failure analysis by extra-prime.
\item Theorem 4: prime 71 in M only among the 26 sporadics.
\end{itemize}

Runtime: <2 seconds. Dependencies: sympy, math.

\section{Status}

\begin{itemize}
\item \textbf{Submission-ready (2026-05-27).} Tier 1.
\item \textbf{Four theorems} with explicit Tier labels (A, A, B, A).
\item \textbf{Tier discipline}: two structural anchors are established --- Theorem 2 (polynomial-vs-factorial dichotomy at rank 24) and Theorem 4 (Ogg's supersingular-prime characterization, anchoring 71 as the genus-0 $X_0(p)^+$ upper bound; subsequently embedded in monstrous moonshine by Conway-Norton 1979 and Borcherds 1992). Theorems 1 and 3 are empirically verified at machine precision with explicit failure analysis.
\item \textbf{No deep moonshine claim}: the strata-prime / supersingular-prime overlap is exact only at the extremes; the intermediate-prime gap is recorded as an open structural question in §7.2.
\end{itemize}

\medskip\hrule\medskip

\textit{--- Sanders \& Gish, 2026-05-27.}

\end{document}